\documentclass[11pt,fleqn]{article}

\usepackage[margin=1in]{geometry}
\usepackage{parskip}
\usepackage{enumitem}
\usepackage{amsmath}    % Better math support
\usepackage{amssymb}    % Optional: additional math symbols

\usepackage{tikz}
\usetikzlibrary{calc}

\begin{document}

%==================================================
% Title
%==================================================

\begin{flushleft}
    {\LARGE \textbf{Álgebra en FEC (I)}}\\[0.5em]
    {\large Ferraris Domingo}\\[0.5em]
    \today
\end{flushleft}

\vspace{0.5em}
\hrule
\vspace{1em}

% ok can you do the same but operating in GF(16) defined by primitive poly x^4+x^3+1, and do the division?:
% (a^8x^3+a^9x^2+a^12x+a^12) / (a^3x^2+a^7x+a^2)
% THEN
% can you give me the latex to show this using tikz?

\[
\begin{tikzpicture}[font=\normalsize]

% Quotient
\node at (5.0,0) {$a^5x+a^{10}$};

% Division symbol and bar
\draw (2.4,-0.35) -- (8.6,-0.35);
\draw (2.4,-0.35) -- (2.4,-5.8);

% Divisor
\node[anchor=east] at (2.35,-0.9)
{$a^3x^2+a^7x+a^2$};

% Dividend
\node[anchor=west] at (2.55,-0.9)
{$a^8x^3+a^9x^2+a^{12}x+a^{12}$};


% First multiplication
\node[anchor=west] at (3.0,-1.7)
{$a^8x^3+a^{12}x^2+a^7x$};

\draw (3.0,-1.95) -- (7.2,-1.95);


% First remainder
\node[anchor=west] at (3.4,-2.6)
{$a^{13}x^2+a^2x+a^{12}$};


% Second multiplication
\node[anchor=west] at (3.4,-3.5)
{$a^{13}x^2+a^2x+a^{12}$};

\draw (3.4,-3.75) -- (7.0,-3.75);


% Final remainder
\node[anchor=west] at (5.0,-4.5)
{$0$};

\end{tikzpicture}
\]

%==================================================
% Summary
%==================================================

\section*{Resumen}

\subsection*{Ejercicio 1 Cambio de representación y suma}
\begin{enumerate}[label=\alph*)]
\item \boxed{\alpha^3+\alpha^9=\alpha}
\item \boxed{\alpha^{11}+\alpha^{13}=\alpha^4}
\item \boxed{\alpha^3+\alpha^9+\alpha^{11}+\alpha^{13}=1}
\item \boxed{\alpha^5+\alpha^5=0}
\end{enumerate}

\subsection*{Ejercicio 2 Producto y división}
\begin{enumerate}[label=\alph*)]
\item \boxed{\alpha^6\cdot\alpha^{11}=\alpha^2}
\item \boxed{\alpha^{12}\cdot\alpha^9=\alpha^6}
\item \boxed{\frac{\alpha^3}{\alpha^{10}}=\alpha^8}
\item \boxed{\frac{1}{\alpha^7}=\alpha^8}
\item \boxed{(\alpha^{14})^3=\alpha^{12}}
\end{enumerate}

\subsection*{Ejercicio 3 Expresiones mixtas}
\begin{enumerate}[label=\alph*)]
\item \boxed{\frac{(\alpha^5+\alpha^9)\cdot\alpha^7}{\alpha^2+\alpha^{12}}=\alpha^6}
\item \boxed{(\alpha^4+1)^2=\alpha^2}
\end{enumerate}

\subsection*{Ejercicio 4 Ecuación lineal}
\boxed{\alpha^5 \cdot x+\alpha^{11}=\alpha^2 \iff x=\alpha^4}

\vspace{1em}
\hrule
\vspace{1em}

%==================================================
% Details
%==================================================

\section*{Desarrollo}

\subsection*{Ejercicio 1 Cambio de representación y suma}
\begin{enumerate}[label=\alph*)]
\item \boxed{\alpha^3+\alpha^9=\alpha}
\begin{align*}
    \alpha^3+\alpha^9
    &=\alpha^3+(\alpha^3+\alpha)\\
    &=(\alpha^3+\alpha^3)+\alpha\\
    &=0+\alpha\\
    &=\alpha
\end{align*}

\item \boxed{\alpha^{11}+\alpha^{13}=\alpha^4}
\begin{align*}
    \alpha^{11}+\alpha^{13}
    &=(\alpha^3+\alpha^2+\alpha)+(\alpha^3+\alpha^2+1)\\
    &=(\alpha^3+\alpha^3)+(\alpha^2+\alpha^2)+\alpha+1\\
    &=0+0+\alpha+1\\
    &=\alpha+1=\alpha^4
\end{align*}

\item \boxed{\alpha^3+\alpha^9+\alpha^{11}+\alpha^{13}=1}\\ \\
Verificar el inciso (c) de dos maneras: sumando los cuatro términos de una vez, y sumando los
resultados obtenidos en (a) y (b).\\
\rule{4cm}{0.4pt}\\\textbf{\textbar \quad 1.}
\begin{align*}
    \alpha^3+\alpha^9+\alpha^{11}+\alpha^{13}
    &=\alpha^3+(\alpha^3+\alpha)+(\alpha^3+\alpha^2+\alpha)+(\alpha^3+\alpha^2+1)\\
    &=(\alpha^3+\alpha^3+\alpha^3+\alpha^3)
    +(\alpha^2+\alpha^2)
    +(\alpha+\alpha)
    +1\\
    &=0+0+0+1\\
    &=1
\end{align*}
\rule{4cm}{0.4pt}\\\textbf{\textbar \quad 2.}
\begin{align*}
    \alpha^3+\alpha^9+\alpha^{11}+\alpha^{13}
    &=\alpha+\alpha^4\\
    &=\alpha+\alpha+1\\
    &=0+1\\
    &=1
\end{align*}

\item \boxed{\alpha^5+\alpha^5=0}
\begin{align*}
    \alpha^5+\alpha^5
    &=(\alpha^2+\alpha)+(\alpha^2+\alpha)\\
    &=(\alpha^2+\alpha^2)+(\alpha+\alpha)\\
    &=0+0=0\\
\end{align*}
\end{enumerate}

\subsection*{Ejercicio 2 Producto y división}
\begin{enumerate}[label=\alph*)]
\item \boxed{\alpha^6\cdot\alpha^{11}=\alpha^2}
\begin{align*}
    \alpha^6\cdot\alpha^{11}
    &=\alpha^{(6+11)\%15}\\
    &=\alpha^{17\%15}\\
    &=\alpha^2
\end{align*}

\item \boxed{\alpha^{12}\cdot\alpha^9=\alpha^6}
\begin{align*}
    \alpha^{12}\cdot\alpha^9
    &=\alpha^{(12+9)\%15}\\
    &=\alpha^{21\%15}\\
    &=\alpha^6
\end{align*}

\item \boxed{\frac{\alpha^3}{\alpha^{10}}=\alpha^8}
\begin{align*}
    \frac{\alpha^3}{\alpha^{10}}
    &=\alpha^3\cdot\alpha^{-10}\\
    &=\alpha^3\cdot\alpha^{(15-10)\%15}\\
    &=\alpha^3\cdot\alpha^5\\
    &=\alpha^8
\end{align*}

\item \boxed{\frac{1}{\alpha^7}=\alpha^8}
\begin{align*}
    \frac{1}{\alpha^7}
    &=\alpha^{-7}\\
    &=\alpha^{(15-7)\%15}\\
    &=\alpha^8
\end{align*}

\item \boxed{(\alpha^{14})^3=\alpha^{12}}
\begin{align*}
    (\alpha^{14})^3
    &=\alpha^{(14\cdot3)\%15}\\
    &=\alpha^{(42)\%15}\\
    &=\alpha^{12}
\end{align*}

\end{enumerate}

\subsection*{Ejercicio 3 Expresiones mixtas}
\begin{enumerate}[label=\alph*)]
\item \boxed{\frac{(\alpha^5+\alpha^9)\cdot\alpha^7}{\alpha^2+\alpha^{12}}=\alpha^6}
\begin{align*}
    \frac{(\alpha^5+\alpha^9)\cdot\alpha^7}{\alpha^2+\alpha^{12}}
    &=\frac{((\alpha^2+\alpha)+(\alpha+\alpha^3))\cdot\alpha^7}{\alpha^2+\alpha^{12}}\\
    &=\frac{(\alpha^3+\alpha^2)\cdot\alpha^7}{\alpha^2+\alpha^{12}}\\
    &=\frac{(\alpha^3+\alpha^2)\cdot\alpha^7}{\alpha^2+(1+\alpha+\alpha^2+\alpha^3)}\\
    &=\frac{(\alpha^3+\alpha^2)\cdot\alpha^7}{1+\alpha+\alpha^3}\\
    &=\frac{(\alpha^3+\alpha^2)\cdot\alpha^7}{\alpha^7}\\
    &=(\alpha^3+\alpha^2)\cdot1\\
    &=\alpha^6
\end{align*}

\item \boxed{(\alpha^4+1)^2=\alpha^2}
\\ \\
Para el inciso (b), resolverlo por dos caminos: primero expresando $\alpha^4+1$ como potencia de $\alpha$ y elevando al cuadrado; luego, multiplicando los polinomios directamente y reduciendo con la
relación $\alpha^4=1+\alpha$. Ambos resultados deben coincidir.

\rule{4cm}{0.4pt}\\\textbf{\textbar \quad 1.}
\begin{align*}
    (\alpha^4+1)^2
    &=((\alpha+1)+1)^2\\
    &=\alpha^2
\end{align*}
\rule{4cm}{0.4pt}\\\textbf{\textbar \quad 2.}
\begin{align*}
    (\alpha^4+1)^2
    &=(\alpha^4+1)(\alpha^4+1)\\
    &=\alpha^4\cdot\alpha^4+\alpha^4+\alpha^4+1\\
    &=\alpha^8+1\\
    &=(\alpha^2+1)+1\\
    &=\alpha^2
\end{align*}
\end{enumerate}

\subsection*{Ejercicio 4 Ecuación lineal}
\boxed{\alpha^5 \cdot x+\alpha^{11}=\alpha^2 \iff x=\alpha^4}
\begin{align*}
    \alpha^5\cdot x+\alpha^{11}&=\alpha^2\\
    \alpha^5\cdot x&=\alpha^2-\alpha^{11}=\alpha^2+\alpha^{11}\\
    x&=\frac{\alpha^2+\alpha^{11}}{\alpha^5}\\
    x&=\frac{\alpha^2+(\alpha^3+\alpha^2+\alpha)}{\alpha^5}\\
    x&=\frac{\alpha^3+\alpha}{\alpha^5}\\
    x&=\frac{\alpha^9}{\alpha^5}\\
    x&=\alpha^9\cdot\alpha^{-5}=\alpha^9\cdot\alpha^{(15-5)\%15}\\
    x&=\alpha^9\cdot\alpha^{10}\\
    x&=\alpha^{19\%15}\\
    x&=\alpha^4
\end{align*}

Verificar el resultado reemplazándolo en la ecuación original.
\begin{align*}
    \alpha^5\cdot(\alpha^4)+\alpha^{11}&\overset{?}{=}\alpha^2\\
    \alpha^9+\alpha^{11}&\overset{?}{=}\alpha^2\\
    (\alpha^3+\alpha)+(\alpha^3+\alpha^2+\alpha)&\overset{?}{=}\alpha^2\\
    0+\alpha^2+0&=\alpha^2
\end{align*}

\vspace{1em}
\hrule
\vspace{1em}

%==================================================
% Conclusion (optional)
%==================================================

% \section*{Conclusion}

% Brief concluding remarks.

\end{document}